Comprendre la conception de la plateforme LaDocumenta.eu, bibliothèque numérique spécialisée en musique, implique d'analyser d'abord le cadre général des bibliothèques numériques, leurs usages et leurs enjeux, pour aborder ensuite le champ plus spécifique des bibliothèques numériques musicales.

\subsection*{Les bibliothèques numériques : cadre général, usages et enjeux}

S'intéresser dans un premier temps au cadre général des bibliothèques numériques à travers la diversité de leurs usages et enjeux constitue un préalable nécessaire pour comprendre comment les collections patrimoniales de ces bibliothèques numériques peuvent se transformer, par la suite, en données structurées exploitables pour la recherche et pour les besoins de la pratique historiquement informée.

\subsubsection*{Les bibliothèques numériques : cadre général, définition et historique}

\begin{quotation}
	Les bibliothèques numériques sont des organisations qui offrent des ressources, y compris en personnel, pour sélectionner, structurer, offrir un accès intellectuel, interpréter, distribuer, conserver des documents sous forme numérique. Une bibliothèque numérique garantit un accès sur la durée aux oeuvres électroniques afin qu'elles soient aisément, et à moindre coût, accessibles par un ou plusieurs publics spécifiques\footnote{Définition de la Digital Library Federation, dans Béquet (Gaëlle), \textit{Trois bibliothèques européennes face à Google.} Publications de l’École nationale des chartes, 2015, \url{https://doi.org/10.4000/books.enc.14063.}}.
\end{quotation}

Aux origines des bibliothèques numériques, citons le projet de biblothèque textuelle numérique de l'université de l'Illinois : le projet Gutenberg amorcé en 1971 par Michael Hart, et les numérisations patrimoniales des bibliothèques dès les années 1980. Ces premières initiatives, associées à l'apparition d'Internet, ont posé les fondations des bibliothèques numériques d'aujourd'hui. 

À l'échelle de la France, la Bibliothèque nationale de France s'est très tôt inscrite dans la volonté créer une bibliothèque numérique avec Gallica en 1997. Comme le rapelle Pierre Carbone, Gallica a été créée pour devenir \enquote{la bibliothèque de l'honnête homme}. D'abord conçue pour numériser des corpus d'oeuvres choisies pour constituer un corpus raisonné des grands textes de la culture française, Gallica a spectaculairement élargi son spectre en numérisant également de la presse. 

Parallèlement aux actions de la BnF, les bibliothèques françaises se sont également dotées d’initiatives portées notamment par le Ministère de l’Enseignement Supérieur et de la Recherche, dans une visée de science ouverte. Il s’agit de portails de valorisation de la littérature scientifique comme Persée- créée en 2005 pour les revues de Sciences Humaines et Sociales- ainsi que des plateformes numériques de  bibliothèques universitaires - comme Medica de la bibliothèque interuniversitaire de Médecine (BIU Santé de Paris-Cité).

\medskip La diversification des plateformes françaises a progressivement nécessité une interopérabilité des fonds, c’est-à-dire une capacité d’échange et de connexion, matérialisée par la création des partenariats de Gallica. 
En effet, Gallica noue des partenariats avec des bibliothèques - comme la Bibliothèque nationale et universitaire de Strasbourg, ou encore la bibliothèque L’Héritage des ponts et chaussées - afin que celles ci puissent numériser leurs documents grâce à une plateforme Gallica marque-blanche. 
Sur le plan technique, cette mise en réseau repose sur l'adoption de standards partagés d'échange de données, comme le protocole OAI-PMH (\textit{Open Archives Initiative Protocol for Metadata Harvesting}). 
En séparant le stockage des fichiers et la diffusion de leurs notices, ce protocole permet à un serveur central (l'agrégateur) de collecter automatiquement les métadonnées de différentes institutions partenaires. Grâce à cette harmonisation (notamment au moyen le format Dublin Core), le chercheur accède à un catalogue national unifié, tandis que chaque établissement conserve la gestion directe de ses collections. L'interopérabilité devient ainsi le pilier de la mutualisation patrimoniale, préfigurant les grands modèles d'agrégation européens. 

Afin de mieux coordonner ces initiatives, un Schéma numérique des bibliothèques a été dessiné en 2010 à la suite duquel une commission Bibliothèques Numériques a suivi ces directives de 2011 à 2014
\footnote{Carbone (Pierre), \textit{Numérique et archivage pérenne : regard du président de la Commission Bibliothèques numériques}, Bulletin des bibliothèques de France (BBF), 2013, n° 5, p. 43-43.}.

À l’échelle européenne, la prise de conscience des enjeux de souveraineté culturelle s’est accélérée en 2004 avec l’annonce du programme Google Print (aujourd’hui connu sous le nom de Google Books). En réaction à cette annonce, six chefs d’Etat et de gouvernement européens - dont la France - ont appelé à la création d’une structure publique commune (2005). 
En 2008, cette idée a abouti à la création d’Europeana après une recommandation de la Commission européenne émise en 2006 sur la \enquote{numérisation et l’accessibilité en ligne du matériel culturel et la conservation numérique}.
Europeana est une bibliothèque numérique européenne dont l’objectif est de mettre à disposition du public un patrimoine européen multilingue illustrant l’éclectisme du patrimoine culturel européen et la richesse de la production scientifique. C’est un portail agrégateur de données, cela veut dire qu’elle ne stocke pas l’ensemble des documents mais moissonne et harmonise leur métadonnées pour unifier les collections des bibliothèques, musées et archives.

\subsubsection*{Usages et enjeux des bibliothèques numériques}

Si \enquote{la numérisation concerne tous les types de documents, pour la plupart dans le domaine public, en fonction de leur intérêt patrimonial, de leur intérêt documentaire (selon les attentes des chercheurs) et de leur état matériel pouvant nécessiter leur préservation\footnote{Carbone (Pierre), \textit{Chapitre X : Les bibliothèques numériques}. Dans : \textit{Les bibliothèques}. Paris: Presses Universitaires de France. Que sais-je ?, 2017.p. 94.}}, la création de bibliothèques numériques dépasse la conservation physique des documents. 
L'évolution des bibliothèques numériques et de leurs cadres institutionnels façonne de nouveux usages pour ces plateformes. 
Dès lors, au-delà des préoccupations initiales de conservation et de diffusion, toujours présentes, la numérisation reconfigure les modes d'accès aux sources et les pratiques de consultation. 
Les documents numérisés, après avoir été des reproductions en format image, sans possibilité de rechercher en mode plein texte, deviennent des données structurées, exploitables enrichies de métadonnées de description pour les contextualiser. 

Si la première génération de bibliothèques numériques a permis la mise à disposition de nombreux fonds patrimoniaux à distance, toutefois, cette dynamique a également permis de souligner leurs enjeux techniques, juridiques et méthodologiques.

Le principal enjeu découle d’une limite des premières campagnes de numérisation. Les fichiers produits y étaient alors numérisés en mode image, sans possibilité de recherche plein texte\footnote{\textit{Ibid.}}. Ne pouvant faire l’objet de campagne d’OCR (Optical Character Recognition) ou d’OMR (Optical Music Recognition), pour la musique, les documents n’étaient pas visibles par les moteurs de recherche. Dès lors, la découvrabilité des sources reposait exclusivement sur la précision des métadonnées du catalogue (auteur, titre, date). Cela compliquant la recherche de sources des chercheurs, qui devaient reproduire le travail préalable à une consultation de sources physiques.

Cette impossibilité d’interroger des corpus de textes, ou de partitions, met en exergue la différence entre la simple mise à disposition de documents à distance et l’environnement de recherche visé lors de la création d’une bibliothèque numérique.L’objectif de ces bibliothèques est de favoriser les réseaux de professionnels de l’information et de chercheurs. 

À ces contraintes matérielles de présentation des documents, s’ajoutent des enjeux de structuration des bibliothèques numériques et de stratégie documentaire. En effet, la numérisation patrimoniale ne peut être neutre. Elle résulte de choix institutionnels guidés, comme l’écrivait Pierre Carbone par \enquote{leur intérêt patrimonial}, \enquote{leur intérêt documentaire} ou encore \enquote{leur état matériel}\footnote{Carbone (Pierre), \textit{Chapitre X : Les bibliothèques numériques}. Dans : \textit{Les bibliothèques}. Paris: Presses Universitaires de France. Que sais-je ?, 2017, p. 94.}.

De ce fait, il est possible de créer des biais d’accès entraînant la surreprésentation de fonds aux dépens d’autres fonds. À ces choix stratégiques et financiers s’ajoutent des contraintes juridiques telles que le cadre strict du droit de la propriété intellectuelle qui cantonne les bibliothèques numériques à la proposition d'œuvres du domaine public. Notons toutefois qu’il existe des exceptions pour les œuvres indisponibles. Le Parlement européen, dans sa directive 2019/790 couvre la numérisation des bibliothèques numériques et prévoit en effet des exceptions au droit d’auteur dans des cas précis et encadrés, notamment pour la conservation du patrimoine culturel et l'exploitation des œuvres indisponibles, afin de favoriser l'accès au savoir.\footnote{\textit{Ibid.}, p. 101-102.}.


Enfin, la gestion de la masse documentaire que représentent les bibliothèques numériques ainsi que la pérennité des documents conservés constitue un autre défi technique et économique des bibliothèques numériques. La conservation d’une telle masse documentaire implique des infrastructures de stockage sécurisées ainsi que des politiques d’archivage pérenne, pour préserver de l’obsolescence tant matérielle que logicielle, ainsi qu’un effort de normalisation des données. 
La transformation des données hétérogènes en données structurées, visant l’interopérabilité  passe par le recours à des modèles de données orientés entité (web sémantique), ou encore au protocole IIIF (\textit{International Image Interoperability Framework}). Ces normes permettent d’éviter l’invisibilité des données et de pouvoir réutiliser des corpus, dans un souci de découvrabilité des ressources numériques d’une institution ou de science ouverte. \footnote{Ces enjeux d’interopérabilité seront approfondis dans le chapitre 5 qui lui sera dédié.}

L'évolution des bibliothèques numériques a transformé la relation aux sources patrimoniales de la communauté scientifique. 
Avec l’accroissement documentaire en ligne, “ la qualité et la visibilité des bibliothèques numériques deviennent essentielles” \footnote {\textit{Ibid}}
Ces plateformes, créées tant pour le grand public que pour les chercheurs se nourrissent des recommandations des usagers, notamment en ce qui concerne leurs pratiques. Ainsi, trois grandes typologies d'usages qui reflètent les mutations techniques de ces plateformes se distinguent:\\

Le premier usage des bibliothèques numériques concerne la consultation numérique distante et immédiate d'un document. Les bibliothèques numériques permettent alors à l'usager, chercheur, étudiant ou grand public de dépasser les limites physiques de consultation : grâce à elles un document non communicable pour des raisons de conservation, ou une impossibilité de se rendre en salle de lecture (dans le cas ou l'ouvrage serait conservé à l'étranger) peuvent être surmontées. Par ailleurs, l'accès à ces documents numériques à haute résolution favorise leur analyse matérielle (codicologie, examen des annotations manuscrites, filigranes) sans dégradation des originaux et avec une préicision accrue permise par le zoom par exemple.\\

Le deuxième usage classique des bibliothèques numériques consiste en l'utilisation des recherches plein texte.
La généralisation des traitements d'OCR et d'OMR modifie structurellement les modalités d'interrogation du document. L'accès ne s'effectue plus seulement par les notices de catalogues et les index d'auteurs mais également en recherchant directement dans le corps du texte ou de la notation musicale. La recherche par mots-clés ou par motifs dans de vastes corpus (collections de presse de Gallica, périodiques scientifiques de Persée) permet notamment de repérer des occurrences rares. \\

Par ailleurs, les humanités numériques et l'approche des \enquote{\textit{Collections as Data}} constituent un autre champ d'application des bibliothèques numériques. Ce modèle consacre le passage du document numérisé au jeu de données structuré, téléchargeable et exploitable par l'intermédiaire d'API. Ces corpus textuels, visuels ou musicaux alimentent désormais des méthodologies d'analyse quantitative : fouille de textes et de données (\textit{text and data mining}), modélisation probabiliste de thématiques (\textit{topic modeling}), cartographie de réseaux d'acteurs, ainsi que l'analyse automatisée de documents visuels et sonores par apprentissage automatique.

La bibliothèque numérique ne se limite plus dès lors à un magasin de bibliothèque en ligne mais devient un réseau de recherche préparant la transition vers l'analyse automatisée de corpus spécialisés.\footnote{Ce développement sur les usages des bibliothèques numériques a utilisé avec profit le cours de culture numérique consacré aux bibliothèques numériques dispensé par M. Canteaut en première année de master. }


\subsection*{La création de bibliothèques numériques en bibliothèque musicale et d'orchestre et les enjeux UX de leur interface en ligne, l'exemple de LaDocumenta.eu}

Après avoir défini le cadre général des bibliothèques numériques, il convient de recentrer l'analyse sur les institutions spécialisées en musique à l'instar des bibliothèques d'orchestre. L'enjeu est de voir comment la création de plateformes comme LaDocumenta.eu et la prise en compte de l'expérience utilisateur (UX) structurent ce patrimoine spécialisé, en le rendant exploitable pour la recherche musicologique et l'interprétation historiquement informée.  

Les bibliothèques numériques spécialisées en musique constituent un type particulier de bibliothèques numériques présentant des contraintes spécifiques, à la fois opérationnelles et techniques. En effet, une bibliothèque d'orchestre ne se limite pas à la conservation de monographies consacrées à la musique comme on pourrait le penser. L'essentiel des fonds se composant plutôt du matériel d'exécution - c'est-à-dire de la partition du chef d'orchestre (le conducteur) et de l'ensemble des partitions individuelles de chaque instrumentiste  (les parties séparées formant le matériel d'orchestre) - de chaque œuvre jouée ou amenée à être jouée par l'orchestre. Le conducteur du chef d’orchestre est un document très riche en informations, dont la lecture peut s’avérer complexe  et pour lequel la synchronisation entre les parties individuelles et la vue d’ensemble est cruciale pour la réussite de l’exécution.

La numérisation et l’organisation de ce matériel d’exécution imposent donc d’adapter la conception des interfaces UX (\textit{User Experience}) en ligne aux spécificités musicales. 

La plateforme LaDocumenta.eu est l’exemple d’interface en ligne  de bibliothèques qui ne peut se contenter d’être un dépôt de PDF en raison des difficultés de prise en main de l’outil. Au contraire, ce type de plateforme présente des enjeux ergonomiques pour répondre aux attentes des usagers en fonction de leur niveau de connaissances de ce milieu spécifique et de leur degré d’autonomie.
LaDocumenta.eu est un projet qui illustre ces enjeux au sein d'une bibliothèque numérique très spécialisée dans le champ de la pratique historiquement informée. 
Rappelons-nous que LaDocumenta.eu est une plateforme collaborative portée par un consortium d'orchestres européen, LaDocumenta. Ce consortium fédère les orchestres \textit{Insula Orchestra}, l'Accademia Bizantina, le Concentus Musicus Wien, le Collegium 1704 et le Balthasar Neumann autour d'un projet de patrimonialisation de la pratique historiquement informée en partenariat avec la recherche académique en musicologie (l'IreMus, le Nikolaus Harnoncourt zentrum) mais aussi dans le champ des humanités numériques pour la construction de leur plateforme singulière : LaDocumenta.eu.

Comme nous l'avons vu, LaDocumenta désigne donc tout autant le consortium d'orchestres que la plateforme en ligne qui en découle, c'est-à-dire LaDocumenta.eu, une bibliothèque musicale numérique qui présente des particularités. 
En effet, LaDocumenta.eu n'est pas une bibliothèque numérique classique qui serait pensée uniquement pour consulter des partitions numérisées. 

Elle possède cette fonctionnalité mais présente en plus des caractéristiques de SIGB (système intégré de gestion de bibliothèque) à l'instar de Koha ou PMB. L'idée de cette plateforme est en effet de combiner une interface utilisateur de médiation de la musique historiquement informée - avec la possibilité de visionner une œuvre numérisée, comme un Gallica de la musique historique - et un système clé en main de gestion de bibliothèque pouvant à terme remplacer un PMB. 
Les orchestres de musique historiquement informée, n'ont pas toujours les moyens pour souscrire à un abonnement PMB et LaDocumenta.eu est pensée pour proposer aux orchestres une interface utilisateur pour la médiation de leur contenu (de type Gallica, très spécialisé en musique) associée à un SIGB. Par exemple, l'orchestre \textit{Insula Orchestra} prévoit de résilier son contrat PMB quand la plateforme sera livrée et parfaitement fonctionnelle. 

LaDocumenta.eu se construit autour de deux éléments majeurs que sont l'ergonomie, pour garantir une expérience utilisateur optimale, dans un souci de développement UX, et la rigueur bibliothéconomique. En effet, si LaDocumenta.eu doit permettre à un musicien novice en bibliothèque de cataloguer une œuvre, elle doit tout de même respecter les standards de bibliothèque pour garantir la reconnaissance de la communauté scientifique, et ne pas décontenancer les authentiques bibliothécaires, spécialisés en musique, qui seront amenés à l'utiliser. 

Pour respecter ces exigences, LaDocumenta.eu s'est inspirée du cen (le centre de ressources  pour l'art choral), lui aussi créé par Laurence Equilbey.
Du modèle du Cen à la version actuelle de LaDocumenta.eu, nous retrouvons l'enjeu majeur des bibliothèques numériques classiques qui est de dépasser le simple site joli, pensé comme une \enquote{vitrine}, pour proposer un site permettant la structuration des données dans une logique d'interopérabilité. 

Dans la première version de la plateforme LaDocumenta.eu, ce critère n'était pas respecté puisque l'interface utilisateur et son affichage, où l'indexation était créée au moyen de tags, prévalaient sur la structuration des ressources. 

Si la plateforme semblait visuellement être une jolie vitrine ergonomique, cette structuration des ressources par tags présentait en réalité de nombreuses limites comme l'absence de traitement des doublons rendant l'exploitation des ressources et leur analyse statistique compliquées, voire impossibles.

Afin de dépasser ces limites, le consortium a réfléchi à une deuxième version de LaDocumenta.eu, dans l'objectif de structurer les ressources de la plateforme pour rendre ses données interopérables.
Il s'agit d'un enjeu important pour intégrer le réseau des musicales européennes et pouvoir imaginer utiliser l'IA pour réaliser certaines tâches.

Mais avant d'utiliser l'IA, il est nécessaire de connaître les usages des utilisateurs. Le livrable rédigé sur les \textit{scenarii} d'utilisation de l'IA dans la plateforme explique dans un premier temps les usages des utilisateurs avant de réfléchir dans un deuxième temps au rôle que pourrait assurer l'IA, pour aider les utilisateurs et contributeurs dans leurs usages.

Pour avoir un ordre d'idée des usages de la plateforme, les développeurs de la plateforme ont réalisé une étude d’usages qu’ils ont soumis aux usagers  lors d’un colloque du consortium. L'objectif était de répartir les usagers dans différentes catégories afin de comprendre leurs attentes et d'adapter ainsi la plateforme aux usages et besoins réels des utilisateurs. 
Comme escompté, des catégories d’utilisateurs se sont distinguées lors de cet exercice. Elles sont au nombre de quatre : 

- En premier lieu, la plateforme répond au besoin d'ingénierie documentaire des professionnels de la conservation ou de la recherche (bibliothécaires, notamment d'orchestre, archivistes et musicologues). Ces utilisateurs, responsables de l'alimentation des fonds, manipulent au quotidien des séries documentaires volumineuses et disparates. 
Pour ce groupe fondamental d'utilisateurs de la plateforme, l'enjeu est de réduire la pénibilité des tâches liées au catalogage. L'idéal serait de proposer des formulaires capables d'analyser des descriptions rédigées en langage naturel. Par exemple : lorsqu'un bibliothécaire écrit \enquote{Documents de travail de l'Enlevèvement au Sérail de Mozart, session de juin 2026, \textit{Insula Orchestra}. Bibliographie.}, le système extrait automatiquement les entités nommées pour pré-remplir les champs normés (titre de la ressource, oeuvre liée, auteur, typologie). Le rôle du professionnel sera alors de contrôler le pré-remplissage de la plateforme. 
En complément, pour plus de visibilité des collections, les professionnels auraient besoin de pouvoir transcrire des annotations manuscrites portées sur les partitions imprimées (analyse graphique associé à MusicXML) ainsi que de pouvoir traduire automatiquement les notices pour lever les barrières linguistiques qui peuvent empêcher la découvrabilité des ressources, c'est-à-dire la visibilité d'un contenu sur le web.

- En second lieu, l'interface doit satisfaire les demandes ciblées liées à la pratique des musiciens interprètes (chefs d'orchestre, instrumentistes, chanteurs, luthiers et ingénieurs du son). Ces interprètes consultent le matériel d'exécution pour résoudre des problématiques concrètes comme par exemple la prononciation, le choix des instruments d'époque ou la disposition des pupitres. Grâce à l'alignement d'identifiants internationaux (issus de la BnF, du RISM, de Wikidata ou de la base MIMO, une base de donnée pour les instruments de musique : \textit{Musical Instrument Museums Online}), l'interface dépasse la simple recherche par mot-clé. L'exemple typique de ces usages est celui d'un luthier cherchant à restaurer un violon baroque selon les standards historiques pour croiser directement la documentation d'un instrument avec des pièces de musée équivalentes et des traités de facture d'époque.

- En troisième lieu, les enseignants, universitaires et chercheurs en musicologie mobilisent la plateforme pour conduire des analyses scientifiques et préserver la mémoire de la pratique historiquement informée (HIP). L'ergonomie de l'interface est ici pensée comme un assistant d'aide à la recherche et à la navigation sémantique dans la masse documentaire. Un chercheur préparant une étude comparative sur l'interprétation de Mozart dispose d'un espace de travail unifié. Il peut s’appuyer sur un agent d'exploration alimenté par un corpus documentaire strict et vérifié (modèle de données souverain et recherche augmentée), lui garantissant une synthèse rigoureuse sans dérives ni fausses informations. Cet environnement permet de mettre en relation des articles académiques, des sources primaires numérisées et des captations de masterclasses à l'échelle européenne. Sur le plan pédagogique, l'interface offre la possibilité de manipuler les données sonores et d'appliquer des filtres acoustiques à des enregistrements d'orchestres modernes, afin d'adapter les timbres et le diapason pour faire entendre aux étudiants le rendu sur instruments d'époque.

- Enfin, LaDocumenta.eu constitue une vitrine de médiation culturelle à destination du grand public, des mélomanes, des critiques et des mécènes. Pour ces usagers désireux de découvrir l'envers du décor d'une création musicale sans maîtriser le solfège ou le catalogage spécialisé, l'expérience utilisateur repose sur la scénarisation et l'adaptation de la restitution. Lorsqu'un amateur effectue une recherche large du type « concert Messe du couronnement \textit{Insula Orchestra} », la plateforme évite de renvoyer une notice d'archive austère. L'interface adapte son niveau de langage et génère un parcours de découverte interactif sous forme de récit visuel (scrollytelling), associant un texte vulgarisé, une iconographie riche et des illustrations audiovisuelles.
La réussite d'une bibliothèque numérique spécialisée en musique comme LaDocumenta.eu réside donc dans la conception d'une interface adaptable. Celle-ci doit concilier la rigueur d'un catalogage patrimonial scientifique avec la fluidité indispensable aux usages artistiques, à la recherche académique et à la démocratisation culturelle.

















